% !TEX root = ../../../main.tex
% Block card: catalogue.cluster -- "Hierarchical cluster" (Part I, the blocks chapter).
\cardfield{What it does}
Groups windows whose measures are alike into a fixed number of clusters, by building a tree of
merges and cutting it.

\cardfield{When to use it}
After a block that gives each window a row of measures:
\blockref{preprocessing.window_matrix}, \blockref{catalogue.event_windows} after Event shape, or
MiniRocket. It is the clustering arm of the first, one-span version of research question 1;
for clustering a pool by shape, use \blockref{catalogue.shape_cluster}.

\begin{plainterms}[Choosing how many clusters]
The tree joins the two most alike groups again and again until one is left; cutting it gives the
clusters. This block always cuts it into exactly \emph{k} clusters, whether or not the windows
really fall into \emph{k} groups: ask for three and you get three, even from pure noise. Two
numbers on the block's page help you judge: the \emph{silhouette} (near 1, clusters well apart;
near 0, they blur into each other) and the \emph{cophenetic correlation} (how faithfully the tree
keeps the real distances).
\end{plainterms}

\cardfield{What goes in, what comes out}
In go windows with their measures. Out comes a cluster number (1 to \emph{k}) for each window,
or $-1$ for a window that had no measures at all.

\cardfield{Settings that matter}
\emph{k}, the number of clusters. \emph{linkage}, how the distance between two groups is
measured: \texttt{ward} (the default) keeps clusters compact; \texttt{average},
\texttt{complete} and \texttt{single} are there because the choice is still open.

\cardfield{Watch out for}
Every window is clustered, test windows included: the train/test roles are ignored, and the
measures are rescaled using all windows. Do not feed it a Window pool directly: the pool's
bookkeeping numbers (channel, start, length) would be clustered as if they were measures.
